\section{When a False Order Costs Time}
\label{sec:condition}

% VALUE: the textbook says a false dependence costs time only on the critical path; on XDNA 2 the useful question is which false orders in compiled designs lie there, and a condition over the platform's resources answers it from the artifacts
% DECISION: the critical-path condition is textbook and cited; the finding is the condition stated over this platform's resources and the designs' cuts, never "false dependences cost time".
Removing a false dependence gains time only where it lies on the critical path, which is the textbook condition~\citep{fields2003interaction}, so the question on XDNA~2 is which false orders in compiled designs lie there, and we answer it with a condition stated over the platform's resources, which the artifacts alone let us check.

\subsection{The Condition}
\label{subsec:clause}

% VALUE: the condition in two clauses: deleting a false order pays only if the work it lets overlap shares no exclusive resource the rewrite cannot duplicate (a buffer can be renamed, a core tile or DMA channel cannot) and no true dependence on the device or through the host also separates it; each clause follows from a platform fact; it names checkable resources and dependences, so it is not the critical-path condition restated
A false order costs time only where deleting it lets two pieces of work overlap that share no exclusive resource the rewrite cannot duplicate, and that no true dependence, on the device or through the host, also separates.
The first clause follows from the platform, because every engine runs its own work in order (Section~\ref{subsec:platform}), so two pieces of work on one core tile or one DMA channel still run in turn once the order between them is deleted.
A buffer is the exception, since a rewrite can rename it into several slots (Section~\ref{subsec:rings}), whereas a core tile or a DMA channel cannot be duplicated.
The second clause holds because a host step that reads one command's output before it issues the next orders the two commands truly, whatever their device streams allow.
Both clauses name resources and dependences that the analysis reads from the artifacts (Section~\ref{sec:order}), so the condition says which false orders can pay before any duration is measured, rather than restating that the ones on the critical path do.
% SOURCE: the condition and its resource wording (vault paper-cosched-decisions-log, "The lead", "The legality control is a designed single-hazard rewrite; the condition's resource wording", 2026-10-24); the census tags each instance by it, and the re-tag under the duplicable-resource wording changes no verdict on any stack (aie-cosched 460383a, rows e857444).

\subsection{Exposure: the Cut and the Shape}
\label{subsec:exposure}

% VALUE: how a design cuts its work into commands and buffers sets how much work meets the condition; per workload we state the exposure M (share of wall time in commands whose configuration holds two or more operators), the bound b inside them, and the ceiling they imply; a design that runs one operator per configuration exposes nothing to reorder inside a command
How a design cuts its work into commands and buffers sets how much of that work meets the condition.
A design that runs one operator per configuration has nothing on disjoint column groups to overlap inside a command, so only its order across commands remains.
We therefore state, per workload, its \emph{exposure} $M$, which is the share of the wall time spent in commands whose configuration holds two or more operators, together with the bound $b$ on the gain inside those commands (Section~\ref{subsec:removal}).
The two imply a ceiling on the end-to-end gain from inside commands of $1 / \bigl((1 - M) + M / b\bigr)$, which separates a design that exposes little work from one whose exposed work cannot overlap.
% SOURCE: exposure metric and ceiling (aie-cosched registered 2a7990a; script 3d11670; rows a3ef131; partb/README).

% VALUE: the shape sets the balance of the exposed work: in a chunked prefill, attention per chunk grows with the chunk's position while the GEMMs' work per chunk is fixed, so the gain from overlapping the groups across chunks is estimated from each group's busy time and the host time the overlap hides; with few chunks the per-command costs and the pipeline's fill and drain outweigh the overlap, so the reorder ties at one chunk and loses at two before it gains; the gain peaks where the groups' work balances and falls where attention outweighs the GEMMs; the split sets the balance, so the best split moves with the shape
The shape of the workload, which for prefill is the prompt length, sets the balance of the exposed work.
In a prefill cut into chunks of 128 tokens, attention's work per chunk grows with the chunk's position while the GEMMs' work per chunk stays fixed, so overlapping the two column groups across chunks can hide the shorter group's time behind the longer one's.
Over a middle range of prompt lengths, the gain is estimated from the two groups' busy times and the host time the overlap hides (Section~\ref{subsec:shape-eval}), and it peaks where the two groups' work balances.
With few chunks, the per-command costs and the time to fill and drain the pipeline of commands outweigh the overlap, so the reorder ties at one chunk and loses at two before it gains.
Past the balance, attention outweighs the GEMMs and the gain falls again.
The split moves this balance, so the split that minimizes TTFT also changes with the shape.
% SOURCE: same-split shape on the development grids, pre-final: one chunk 0.976-1.006, two chunks 0.945-0.972, four chunks 1.040-1.076 on all five LLMs, 1.13-1.18x at 1,024, 1.82-2.17x at 8-16K, falling at 32K (vault decisions log, grids of 2026-10-23/24; issue check 8d2eac19). The estimate with its host term against the event model and the measured cells: pass-0 group terms (registered 3fab31d) and the synthetic sweep (abaa3d9).
% DECISION: the shape rule is an estimate for a middle range of lengths, never a bound: a plain cap of 2x is contradicted by Llama-3.2-1B's 2.07-2.17x, because it omits the host time the composer hides; at few chunks only the event model prices the per-command costs, and at long prompts where the groups balance, overlap within one command stays short of perfect (Gemma3-4B at 16-32K), which only the event model's per-command composition captures (group-term fit 6c5b666).

% VALUE: the consequence: the order is fixed in the artifact while the exposed work changes with the shape, so the order that pays changes with the shape; choosing it needs durations, which no XDNA 2 compile has, and no compile sees the host's order across commands; so the order and the split must be priced per shape
The order is fixed in the artifact while the work it exposes changes with the shape, so the order that pays changes with the shape, and because choosing it needs durations that no compile has, OPCORD prices the order and the split per shape from timings of each operator run alone (Section~\ref{sec:pricing}).

\subsection{Where the Stacks Meet It}
\label{subsec:observed}

% VALUE: across the 112 designs we read on three stacks, the condition holds in one design family, MLIR-AIR's chunked fused prefill, where the gains are measured; FastFlowLM's flash prefill exposes most of its time in multi-operator commands, yet nothing inside them can overlap and its host steps truly separate most commands; its decode pairs share an engine; IRON runs one operator per configuration; no non-LLM design meets it; so we evaluate every design in the census that meets the condition
Across the 112 designs that we read on three stacks, the condition holds in one family of designs, which is where we measure the gains.
MLIR-AIR's chunked fused prefill issues one command per operator and chunk, with attention and the GEMMs on disjoint column groups, so its column groups take turns on work that meets both clauses (Section~\ref{subsec:losses}).
FastFlowLM's flash prefill, taken as released, spends 82--87\% of its time in commands whose configuration holds several operators, yet the false orders inside them bound at $1.000\times$, because its streams already keep every column busy at once, so deleting those orders lets no further work overlap.
Between its commands, most waits are true, since the host reads each command's output before it issues the next, and the false ones bound below $1.05\times$ on measured terms on every request we price.
Its decode runs one command per step, whose false pairs all share a core tile or a DMA channel, and bounds at $1.000$--$1.009\times$.
IRON's designs run one operator per configuration, and each of the non-LLM designs we read runs as one command whose stages already overlap through the array's own dataflow, so neither meets the condition.
We therefore evaluate every design in the census that meets the condition, and a synthetic LLM on the same engine varies the balance and the number of chunks directly (Section~\ref{subsec:shape-eval}).
% SOURCE: census Part A 112 designs (aie-cosched 7827f85); FastFlowLM exposure, scored request, M 0.82-0.87 (flash), b over the legal rewrite space 1.000 under the task form (terms4, registered 67d54fe, rows 3c70af2; exposure 2a7990a / a3ef131); 211 of 312 waits true at E2B 1,024 (census SUMMARY); measured-only joint below 1.05 on every scored request (terms3/terms4); decode boundary 1.000-1.009, inside a runlist 1.000 (decode2 rows e679463). IRON composition predicted 1.000 for every design (aie-cosched a73ba38 / dbdbda0); non-LLM designs: each one command (pass-0 task 6, 2026-10-24).
% SOURCE: FastFlowLM's only crossing over its legal rewrite space is the dispatch a merge of two commands saves at the 128-token bucket, 1.075-1.093x as an upper bound at the measured floor F = 195.0 us (terms5, aie-cosched e8232aa); the merges are measured in Section~\ref{subsec:census}.
% DECISION: FastFlowLM, taken as released, is an observation of the condition, never a like-for-like comparison and never evidence that the cost is avoidable; no claim about how its order was chosen.
% SOURCE: IRON in-command pricing, 35 at 1.000 = the population's 32 IRON streams plus 3 non-LLM designs (mobilenet, resnet, bottleneck) (aie-cosched 1a87bc6; membership adf0e2c).
% COMMENT: the FastFlowLM merge replay (stage 1) is measured in Section~\ref{subsec:census}.
